% ---------------------------------------------------------------
% Solution representation and training objective.
% Replaces the two-network / SIREN passage.
% ---------------------------------------------------------------

Because the Robin problem~\eqref{eq:pde}--\eqref{eq:robin} is posed on $\Omega^{-}$
alone, the solution and its gradient are continuous across $\Gamma$ and no jump
conditions arise. A single network therefore suffices, in contrast to
jump-condition formulations which require separate representations either side of the
interface. We write
\begin{equation}
    u_{\theta} = \mathcal{N}(\mathbf{x};\, \theta) : \mathbb{R}^{3} \to \mathbb{R}.
    \label{eq:network}
\end{equation}

$\mathcal{N}$ is a fully connected feedforward network with two hidden layers of $100$
units, $\tanh$ activations, and a single linear output neuron, giving $10{,}601$
parameters. It takes the raw coordinate $\mathbf{x}$ as input; no positional or
spectral encoding is applied. Hidden weights are drawn from a truncated normal
distribution with zero mean and standard deviation $0.1$, and biases are initialised
to zero. A smooth nonlinearity is essential: piecewise-linear activations such as
ReLU have vanishing curvature almost everywhere and are unsuitable for representing
solutions of second-order differential equations, whose discrete operator differences
the network across neighbouring points.

Note that $\mathcal{N}$ is defined on all of $\mathbb{R}^{3}$, not on $\Omega^{-}$
alone. The finite-volume stencil at a cell adjacent to $\Gamma$ samples the solution
at $\mathbf{x} \pm \Delta x\, \mathbf{e}_{i}$, and for cells the interface crosses
some of these neighbours lie outside $\Omega^{-}$. Evaluating $\mathcal{N}$ there
supplies a smooth extension of the interior solution across $\Gamma$ at no additional
cost and with no explicit extrapolation. Equations are posed only at points of
$\Omega^{-}$; the representation is not restricted to it.

The network is trained by minimising the residual of the discrete system over the
interior grid points $I = \{\, (i,j,k) : \varphi(\mathbf{x}_{ijk}) \leq 0 \,\}$:
\begin{equation}
    \mathcal{L}(\theta)
    = \frac{1}{|I|} \sum_{(i,j,k) \in I}
      \frac{1}{2}
      \left[
        \frac{\bigl(A\, u_{\theta}\bigr)_{ijk} - b_{ijk}}{d_{ijk}}
      \right]^{2},
    \label{eq:loss}
\end{equation}
where $A$ and $b$ are the finite-volume operator and source assembled from
\eqref{eq:pde}--\eqref{eq:robin}, and $d_{ijk}$ is the diagonal entry of row $ijk$.
The division by $d_{ijk}$ is a Jacobi preconditioning of the residual. Its purpose is
not to accelerate a linear solve -- none is performed -- but to equalise the
contribution of cells of differing size: a cell that the interface cuts may retain a
small fraction of its nominal volume, and without this scaling its equation would
enter~\eqref{eq:loss} with a correspondingly reduced weight, precisely at the cells
where the boundary condition is imposed.

No exact solution enters~\eqref{eq:loss}. The only data are $f$, $g$, $\mu$, $k$,
$\alpha$ and $\varphi$; the values of $u_{\theta}$ at stencil neighbours are supplied
by the network itself, so the discrete system is satisfied in a least-squares sense
over the parameters $\theta$ rather than solved for nodal unknowns.
